% GENERATED FILE. Edit canonical lessons, then run npm run generate:books.
% canonical-source-hash: fnv1a64:1690b64c97d06a8f
% canonical-lessons: UR-C136-pyasa, UR-C136-taza, UR-C136-kachcha, UR-C136-pakka, UR-C136-gulabi

\chapter{Thirsty, Fresh, Raw, Ripe, Pink}
\label{ch:ur-thirsty-fresh-raw-ripe-pink}
\hlchaptermodality{136}

\begin{chapteropening}
\textbf{By the end of this chapter:} \emph{I can say thirsty, fresh, raw, ripe and pink.}

Complete the last lesson of chapter 136: I can say thirsty, fresh, raw, ripe and pink.
\end{chapteropening}

\section[\texorpdfstring{\ur{پیاسا} (pyāsā)}{پیاسا (pyāsā)}]{\ur{پیاسا} (pyāsā) — thirsty}

\label{lesson:UR-C136-pyasa}



\begin{quote}
Before the new one: say the Urdu for tired, then the Urdu for hungry.
\end{quote}

\subsection*{You'll want to know: \ur{پیاسا}}
\textbf{\ur{پیاسا}} — \emph{pyāsā} — \textquotedblleft{}thirsty\textquotedblright{}.

\begin{cousinweb}
\textbf{\ur{پیاسا}} (\emph{pyāsā}) is a tadbhava, a word inherited from Sanskrit and reshaped by sound change. It comes from \emph{*pipāsaka}, \textquotedblleft{}thirsty\textquotedblright{}, a form CDIAL reconstructs beside Sanskrit \emph{pipāsita} (CDIAL 8200).
\end{cousinweb}

\begin{grammarlens}[title={What you've built}]
One more word for this chapter.
\end{grammarlens}

\subsection*{Your turn}
\begin{itemize}
\raggedright
  \item \emph{Say it:} \emph{pyāsā}
  \item \emph{Say it:} \emph{pyāsā}, once more
  \item \emph{Say it:} say \emph{pyāsā}
  \item \emph{Recall:} say the Urdu for tired, then the Urdu for hungry, then say \emph{pyāsā} again
\end{itemize}

\begin{tcolorbox}[breakable,colback=teal!4,colframe=teal!35!black,title={Before you move on}]
What does \ur{پیاسا} mean? (\textquotedblleft{}Thirsty\textquotedblright{}.) Say it once more.
\end{tcolorbox}
\section[\texorpdfstring{\ur{تازہ} (tāzā)}{تازہ (tāzā)}]{\ur{تازہ} (tāzā) — fresh}

\label{lesson:UR-C136-taza}



\begin{quote}
Before the new one: say the Urdu for hungry, then the Urdu for thirsty.
\end{quote}

\subsection*{You'll want to know: \ur{تازہ}}
\textbf{\ur{تازہ}} — \emph{tāzā} — \textquotedblleft{}fresh\textquotedblright{}.

\begin{cousinweb}
\textbf{\ur{تازہ}} (\emph{tāzā}) is a loan, usually traced to Persian \emph{tāza}, \textquotedblleft{}fresh\textquotedblright{}.
\end{cousinweb}

\begin{grammarlens}[title={What you've built}]
One more word for this chapter.
\end{grammarlens}

\subsection*{Your turn}
\begin{itemize}
\raggedright
  \item \emph{Say it:} \emph{tāzā}
  \item \emph{Say it:} \emph{tāzā}, once more
  \item \emph{Say it:} say \emph{tāzā}
  \item \emph{Recall:} say the Urdu for hungry, then the Urdu for thirsty, then say \emph{tāzā} again
\end{itemize}

\begin{tcolorbox}[breakable,colback=teal!4,colframe=teal!35!black,title={Before you move on}]
What does \ur{تازہ} mean? (\textquotedblleft{}Fresh\textquotedblright{}.) Say it once more.
\end{tcolorbox}
\section[\texorpdfstring{\ur{کچا} (kachchā)}{کچا (kachchā)}]{\ur{کچا} (kachchā) — raw, unripe}

\label{lesson:UR-C136-kachcha}



\begin{quote}
Before the new one: say the Urdu for thirsty, then the Urdu for fresh.
\end{quote}

\subsection*{You'll want to know: \ur{کچا}}
\textbf{\ur{کچا}} — \emph{kachchā} — \textquotedblleft{}raw, unripe\textquotedblright{}.

\begin{cousinweb}
\textbf{\ur{کچا}} (\emph{kachchā}) goes back to a reconstructed \emph{*kacca}, \textquotedblleft{}raw, unripe\textquotedblright{} (CDIAL 2613). CDIAL lists Hindi \emph{kācā} and marks \emph{kaccā} as taken over from Punjabi.
\end{cousinweb}

\begin{grammarlens}[title={What you've built}]
One more word for this chapter.
\end{grammarlens}

\subsection*{Your turn}
\begin{itemize}
\raggedright
  \item \emph{Say it:} \emph{kachchā}
  \item \emph{Say it:} \emph{kachchā}, once more
  \item \emph{Say it:} say \emph{kachchā}
  \item \emph{Recall:} say the Urdu for thirsty, then the Urdu for fresh, then say \emph{kachchā} again
\end{itemize}

\begin{tcolorbox}[breakable,colback=teal!4,colframe=teal!35!black,title={Before you move on}]
What does \ur{کچا} mean? (\textquotedblleft{}Raw, unripe\textquotedblright{}.) Say it once more.
\end{tcolorbox}
\section[\texorpdfstring{\ur{پکا} (pakkā)}{پکا (pakkā)}]{\ur{پکا} (pakkā) — ripe, firm}

\label{lesson:UR-C136-pakka}



\begin{quote}
Before the new one: say the Urdu for fresh, then the Urdu for raw.
\end{quote}

\subsection*{You'll want to know: \ur{پکا}}
\textbf{\ur{پکا}} — \emph{pakkā} — \textquotedblleft{}ripe, firm\textquotedblright{}.

\begin{cousinweb}
\textbf{\ur{پکا}} (\emph{pakkā}) goes back to Sanskrit \emph{pakva}, \textquotedblleft{}cooked, ripe\textquotedblright{} (CDIAL 7621). CDIAL marks the Hindi form as taken over from Punjabi; \textquotedblleft{}firm, sure\textquotedblright{} grew out of \textquotedblleft{}ripe, mature\textquotedblright{}.
\end{cousinweb}

\begin{grammarlens}[title={What you've built}]
One more word for this chapter.
\end{grammarlens}

\subsection*{Your turn}
\begin{itemize}
\raggedright
  \item \emph{Say it:} \emph{pakkā}
  \item \emph{Say it:} \emph{pakkā}, once more
  \item \emph{Say it:} say \emph{pakkā}
  \item \emph{Recall:} say the Urdu for fresh, then the Urdu for raw, then say \emph{pakkā} again
\end{itemize}

\begin{tcolorbox}[breakable,colback=teal!4,colframe=teal!35!black,title={Before you move on}]
What does \ur{پکا} mean? (\textquotedblleft{}Ripe, firm\textquotedblright{}.) Say it once more.
\end{tcolorbox}
\section[\texorpdfstring{\ur{گلابی} (gulābī)}{گلابی (gulābī)}]{\ur{گلابی} (gulābī) — pink}

\label{lesson:UR-C136-gulabi}



\begin{quote}
Before the new one: say the Urdu for raw, then the Urdu for ripe.
\end{quote}

\subsection*{You'll want to know: \ur{گلابی}}
\textbf{\ur{گلابی}} — \emph{gulābī} — \textquotedblleft{}pink\textquotedblright{}.

\begin{cousinweb}
\textbf{\ur{گلابی}} (\emph{gulābī}) is a loan, usually traced to Persian \emph{gulābī}, \textquotedblleft{}rose-coloured\textquotedblright{}. It is built on \emph{gulāb}, \textquotedblleft{}rose-water\textquotedblright{}, from \emph{gul}, \textquotedblleft{}rose\textquotedblright{}, and \emph{āb}, \textquotedblleft{}water\textquotedblright{}.
\end{cousinweb}

\begin{grammarlens}[title={What you've built}]
One more word for this chapter.
\end{grammarlens}

\subsection*{Your turn}
\begin{itemize}
\raggedright
  \item \emph{Say it:} \emph{gulābī}
  \item \emph{Say it:} \emph{gulābī}, once more
  \item \emph{Say it:} say \emph{gulābī}
  \item \emph{Recall:} say the Urdu for raw, then the Urdu for ripe, then say \emph{gulābī} again
\end{itemize}

\begin{tcolorbox}[breakable,colback=teal!4,colframe=teal!35!black,title={Before you move on}]
What does \ur{گلابی} mean? (\textquotedblleft{}Pink\textquotedblright{}.) Say it once more.
\end{tcolorbox}
